% \iffalse meta-comment
%
% hit-final-cover.dtx 是封面（封一）
%
% \fi
%
% \section{封面}
%
% \changes{v3.2a}{2026/08/24}{封面与内封拆成 \texttt{cover} 与 \texttt{titlepage} 两个模块；
%   封面改照 Word 范例的段落序列排，本硕博共用一套，空行组合按「校名不动、姓名居中」
%   加权择优，排不下时按六档结构让步；本科三校区统一，非本部本科那几套版式与具名间距
%   删掉；学生类别按 \option{degree-type} 取，新版面研究生都排、本科不排；\cs{makecover} 记下「前置已排」，老写法
%   不再在文末兜底补排；去掉 \cs{end}\{tabular\} 后的词间空格（\issue[294]）}
%
% \emph{这个模块只管封面一。} 内封（扉页）挪到 \file{hit-final-titlepage.dtx}。
% 名字也一并改正：从前叫「第一张封面」「第二张封面」，实际上前者是封面、后者是
% 内封，英文里本来就是 cover 与 titlepage 两个词。
%
%   \cs{hit@first@title@page}          \ttfamily→\rmfamily \cs{hit@cover}
%   \cs{hit@first@title@page@postdoc}  \ttfamily→\rmfamily \cs{hit@cover@postdoc}
%   \cs{hit@first@title@page@other}    \ttfamily→\rmfamily \cs{hit@cover@not@postdoc}
%   \cs{hit@second@title@page}         \ttfamily→\rmfamily \cs{hit@titlepage}
%   \cs{hit@second@title@page@postdoc} \ttfamily→\rmfamily \cs{hit@titlepage@postdoc}
%   \cs{hit@second@title@page@bachelor}\ttfamily→\rmfamily \cs{hit@titlepage@bachelor}
%   \cs{hit@second@title@page@other}   \ttfamily→\rmfamily \cs{hit@titlepage@other}
%   \cs{hit@english@cover}             \ttfamily→\rmfamily \cs{hit@titlepage@en}
%   \cs{hit@cover@line}                \ttfamily→\rmfamily \cs{hit@titlepage@line}
%
% 拆分是纯搬运：四个变体（博士本部、本科本部、博后威海、博士本部英文，共 180 页）
% 逐像素与拆分前相同。
%
%    \begin{macrocode}
%<*final-cover>
\bool_if:NT \g__hit_final_bool
  {
%    \end{macrocode}
%
% 封面标题定宽两端对齐用的 \env{CJKfilltwosides} 由 \pkg{xeCJKfntef} 提供。
% 它还能避开 \cs{makebox}\oarg{width}\oarg{s} 可能产生的 underfull box。
%    \begin{macrocode}
\skip_new:N \hit@title@width
\NewDocumentCommand \hit@put@title { O{\hit@title@width} m }
  { \hit@spread@to:nn {#1} {#2} }
%    \end{macrocode}
%
% 封面上那行学生类别排什么。学术与专业各有一条文本，按 \option{degree-type} 选；
% 用户在 \option{student-type} 里写了东西就以他写的为准。
%
% 判空要跟一个建法完全相同的空宏比，不能用 \cs{tl\_if\_empty:NTF}。元信息字段是
% \cs{cs\_gset\_nopar:cpn} 建的，不带 \cs{long}，而 \cs{c\_empty\_tl} 的
% \cs{long} 与否不由这边决定，比意义会错。同样的坑在 \cs{hit@after@number:nnn}
% 那里踩过一次，症状是本该走默认值的地方静悄悄地什么都不排。
%    \begin{macrocode}

%    \end{macrocode}
%
% 这三个是“字段名加一条下划线”那种封面行型，本科三校区统一之后没有调用者了。
% 内封不画下划线，走的是下面的 \cs{hit@titlepage@line}。
%
% \changes{v3.1d}{2025/03/03}{将哈尔滨本科的第一封面替换为 \cmd{\hit@cover@not@postdoc}。\issue[172]}
% 将哈尔滨本科的第一封面替换为 \cmd{\hit@cover@not@postdoc}。
%
% 第一张封面有两套版式：博后一套，其余那一套是 \cs{hit@cover@not@postdoc}。
% 名字写成取反，不另起一个正面的名字，跟条件里的 \option{layout}\texttt{=!v2-layout}
% 与 \cs{hit@if@not@shenzhen} 是同一条规矩：默认那一档没有自己的身份，它就是
% 「不是那一档」，另起一个名字（\texttt{universal}、\texttt{other}）只会让人以为
% 它是并列的第二种东西。它也不真 universal，博后不在里面。
%
% 等博后的封面搬到新版面，这两套并成一套，取反这个名字连同
% \cs{hit@cover@postdoc} 一并删掉。
% 本科三个校区已经统一，不再按校区分。
%    \begin{macrocode}
\cs_new_protected:Npn \hit@cover
  {
    \bool_if:NTF \g__hit_postdoc_bool
      { \hit@cover@postdoc }
      { \hit@cover@not@postdoc }
  }

%    \end{macrocode}
%
% 添加博后封皮选项
% \changes{v3.0.12}{2020/11/02}{添加博后封皮选项}
%    \begin{macrocode}
\cs_new_protected:Npn \hit@cover@postdoc
  {
~
 {\songti\sihao\cs_set:Npn \arraystretch {2}\ %
 \begin{tabular}{@{}r@{}l@{\hspace{7\ccwd}}r@{}l@{\hspace{3em}}r}~
 \hit@term@cover@classified@index@postdoc&\underline{\makebox[6\ccwd]{\hit@info@classified@index@national}}&\hit@term@cover@secrecy&\underline{\makebox[6\ccwd]{\hit@info@secret@level}}&\\~
 \hit@term@cover@udc@postdoc&\underline{\makebox[6\ccwd]{\hit@info@classified@index@international}}&\hit@term@cover@number@postdoc&\underline{\makebox[6\ccwd]{\hit@info@report@number}}&\\~
 \end{tabular}\cs_set:Npn \arraystretch {1}}~
 \parbox[t][5cm][b]{\textwidth}{\begin{center}\heiti\hit@fontsize@x{xiaoer}{2}\hit@term@cover@university@zh\par
 \hit@spread@by:nn{0.58\ccwd}{\hit@term@cover@report}\end{center}}~
 \parbox[b][4cm][b]{\textwidth}{~
 \songti\hit@fontsize@x{xiaosan}{2}\cs_set:Npn \arraystretch {1.5}~
 \begin{center}~
 \begin{tabular}{c}~
 \makebox[9cm]{\hit@info@title@line@one@zh}\\~
 \hline
 \makebox[9cm]{\hit@info@title@line@two@zh}\\~
 \hline
 \end{tabular}\cs_set:Npn \arraystretch {1}
 \end{center}}~
 \begin{center}~
 \songti\hit@fontsize@x{sihao}{2}\centering\hit@info@author@zh
 \end{center}~
 \parbox[b][4.5cm][b]{\textwidth}{~
 \songti\hit@fontsize@x{sihao}{2}\cs_set:Npn \arraystretch {1.5}~
 \begin{center}~
 \begin{tabular}{rl}~
 \hit@term@cover@work@finish@date@postdoc &~\underline{\makebox[6cm]{\hit@info@work@finish@date@zh}}\\~
 \hit@term@cover@report@submit@date@postdoc &~\underline{\makebox[6cm]{\hit@info@report@submit@date@zh\hfill}}\\~
 \end{tabular}\cs_set:Npn \arraystretch {1}
 \end{center}~
 }~
 \vfill
 {\songti\hit@fontsize@x{sihao}{2}~
 \begin{center}~
 \hit@term@cover@university@zh%
 （\hit@term@cover@campus）%
 \end{center}~
 \begin{center}~
 \hit@info@date@zh
 \end{center}}~
  }



%    \end{macrocode}
%
% 添加博后封皮
% \changes{v3.0.0}{2020/05/21}{添加博后封皮}
%
% 本科三个校区已经统一，这里的两套本科封面连同它们专用的
% \cs{hit@cover@title@weihai} 与 \file{hit-weihai-bachelor-title.eps} 一并删了。
%
% 封面顶上那行叫法各变体不同，收在这里；纵向位置不再靠具名间距硬凑，见下面的
% 段落序列。
%    \begin{macrocode}
\cs_set_nopar:Npn \hit@cover@document@name { \hit@term@cover@degree@level@zh \hit@term@cover@document@type }
\hit@if@option@TF{bachelor}
  {
    \cs_set_nopar:Npn \hit@cover@document@name
      { \hit@term@cover@document@type@level \hit@term@cover@document@type }
  }
  { }

%    \end{macrocode}
%
% \subsection{封面一}
%
% 不照旧版的 \cs{parbox} 加 \cs{vspace} 排，直接照 Word 范例那一页的段落序列
% 排——本科取本科书写范例的第二页，研究生取研究生书写范例的第二页。范例那一页
% 带批注，每个元素下面跟着「↑」与一行说明；真封面不是把批注段删掉，是把批注的
% 字删掉，段落还在、高度不变，所以每处批注都换成同高的空段。
%
% 本硕博走同一套序列，研究生只在标签之后多插一段学位类别加它的两个空段。
%
% 空段的高按 \texttt{w:rFonts} 的 ascii 槽算，不是 eastAsia 槽。页顶那三段与
% 校名前那六段都没有 \texttt{rPr}，ascii 继承 Times New Roman，一个 13.80；
% 批注段位置上那几个显式写了宋体，一个 15.5625。
%
% 这一页只有中文题目那一段贴网格，别的段都显式写着 \texttt{snapToGrid=0}。
%
%    \begin{macrocode}
\cs_new_protected:Npn \hit@cover@blank@zh:n #1
  {
    \int_compare:nNnT {#1} > { 0 }
      { \hitenter*[lines=#1,font={size=xiaosi,asian-text-font=SimSun},paragraph={spacing={line-spacing=multiple,at=1}}] }
  }
\cs_new_protected:Npn \hit@cover@blank@en:n #1
  {
    \int_compare:nNnT {#1} > { 0 }
      { \hitenter*[lines=#1,font={size=xiaosi,latin-text-font=Times New Roman},paragraph={spacing={line-spacing=multiple,at=1}}] }
  }
%    \end{macrocode}
%
% 姓名上面那一段是二号、1.25 倍，ascii 写着宋体，按汉字算，35.664。校名前那一段
% 是小二、1.25 倍，eastAsia 写着黑体但 ascii 空着，继承 Times，按西文算，25.875
% ——按黑体算成 29.18 是错的。两段都不是单倍空行，不参与增减，只有让步阶梯最后
% 那几级才动：让一级把 1.25 改单倍，再让就整段删掉，姓名上面那个先删。
%    \begin{macrocode}
\cs_new_protected:Npn \hit@cover@fixed@author:n #1
  {
    \int_compare:nNnT {#1} < { 2 }
      {
        \int_compare:nNnTF {#1} = { 0 }
          { \hitenter*[lines=1,font={size=erhao,asian-text-font=SimSun},paragraph={spacing={line-spacing=multiple,at=1.25}}] }
          { \hitenter*[lines=1,font={size=erhao,asian-text-font=SimSun},paragraph={spacing={line-spacing=multiple,at=1}}] }
      }
  }
\cs_new_protected:Npn \hit@cover@fixed@inst:n #1
  {
    \int_compare:nNnT {#1} < { 3 }
      {
        \int_compare:nNnTF {#1} = { 0 }
          { \hitenter*[lines=1,font={size=xiaoer,latin-text-font=Times New Roman},paragraph={spacing={line-spacing=multiple,at=1.25}}] }
          { \hitenter*[lines=1,font={size=xiaoer,latin-text-font=Times New Roman},paragraph={spacing={line-spacing=multiple,at=1}}] }
      }
  }
%    \end{macrocode}
%
% \emph{新版面下学生类别那一行研究生一律排，本科不排}，照 iota 的
% \texttt{degree-type-for}：研究生不写就是学术学位。范例封面上那一行是
% “（学术/专业学位论文）”二选一的占位，哈尔滨博士也有。旧版面照 v3.1e
% 的判据（哈尔滨全日制博士不标），那一档钉死不动。三态判的是一个表达式，
% 要能展开，所以先把版面判完落进一个布尔量。
%    \begin{macrocode}
\bool_new:N \l__hit_cover_degree_type_bool
\cs_new_protected:Npn \__hit_cover_degree_type_default:
  {
    \__hit_layout_msword_active:TF
      { \bool_set:Nn \l__hit_cover_degree_type_bool { ! \g__hit_bachelor_bool } }
      {
        \hit@if@show@degree@type@TF
          { \bool_set_true:N \l__hit_cover_degree_type_bool }
          { \bool_set_false:N \l__hit_cover_degree_type_bool }
      }
  }
%    \end{macrocode}
%
% 页顶到英文题目为止。\meta{挤占} 是挤掉的中英题目之间的空行数。
%    \begin{macrocode}
\cs_new_protected:Npn \hit@cover@upto@en:nn #1#2
  {
    \__hit_cover_degree_type_default:
    \hit@cover@blank@en:n { 3 }
    \docxlike_format_par:nn { cover-label } { \hit@cover@document@name }
    \hit@cover@blank@zh:n { 2 }
    \__hit_tristate_if:NnT \g__hit_cover_degree_type_tl
      { \l__hit_cover_degree_type_bool }
      {
        \docxlike_format_par:nn { cover-type }
          {
            \hit@term@cover@brace@left@zh
            \hit@term@degree@type@zh
            \hit@term@cover@brace@right@zh
          }
        \hit@cover@blank@zh:n { 2 }
      }
    \docxlike_format_par:nn { cover-title } { \hit@info@title@cover@zh }
    \hit@if@subtitle@T
      {
        \docxlike_format_par:nn { cover-subtitle }
          { ——\hit@info@title@subtitle@zh }
      }
    \int_compare:nNnT {#2} < { 2 } { \hit@cover@blank@zh:n { 1 } }
    \int_compare:nNnT {#2} < { 1 } { \hit@cover@blank@zh:n { 1 } }
    \docxlike_format_par:nnn { cover-title-en } { font = { size = #1 } }
      {
        \hit@info@title@cover@en
        \hit@if@subtitle@T
          { \hit@term@cover@title@separator@en \hit@info@title@subtitle@en }
      }
  }
\cs_new_protected:Npn \hit@cover@author@line:
  { \docxlike_format_par:nn { cover-author } { \hit@info@author@zh } }
\cs_new_protected:Npn \hit@cover@inst@lines:
  {
    \docxlike_format_par:nn { cover-institution } { \hit@term@cover@university@zh }
    \docxlike_format_par:nn { cover-date } { \hit@info@cover@date@zh }
  }
%    \end{macrocode}
%
% \subsection{封面一的空行组合}
%
% 姓名与校名那两行的基线\emph{不钉死}，靠增减空段维持在参考位置附近：题目一长
% 就少发一个空段把它提回来，一短就多发一个。空段有最小粒度，落点只能落在那个
% 粒度上，所以是「大体不变」。参考位置实测自 Word 范例第二页。
%
% 楷体\_GB2312 一律按中易楷体算。研究生那份 docx 的封面写的是楷体\_GB2312，
% 与中易楷体的度量本来就不一样，要分开还得单独标定一份系数，眼下两者同一档。
%    \begin{macrocode}
\dim_new:N \l__hit_cover_top_dim
\dim_new:N \l__hit_cover_width_dim
\dim_new:N \l__hit_cover_height_dim
\dim_new:N \l__hit_cover_bc_dim
\dim_new:N \l__hit_cover_be_dim
\dim_new:N \l__hit_cover_ha_dim
\dim_new:N \l__hit_cover_hi_dim
\dim_new:N \l__hit_cover_iasc_dim
\dim_new:N \l__hit_cover_inst_target_dim
\box_new:N \l__hit_cover_box
\int_step_inline:nnn { 0 } { 3 }
  {
    \dim_new:c { l__hit_cover_ba_#1_dim }
    \dim_new:c { l__hit_cover_bi_#1_dim }
  }
\int_step_inline:nnn { 0 } { 2 }
  { \dim_new:c { l__hit_cover_tu_#1_dim } }
%    \end{macrocode}
%
% 名字里带数字的那几个量只能用 \texttt{:c} 变体取——\pkg{expl3} 里数字不算字母，
% 直接写 \cs{l\_\_hit\_cover\_tu\_0\_dim} 会被拆成 \cs{l\_\_hit\_cover\_tu\_}
% 加上「0\_dim」三个字符。
%
% 量一块的高。给版心宽是必须的：不给宽的话英文题目不折行，量出来只有一行，
% 而那正是要判的那件事。
%    \begin{macrocode}
\cs_new_protected:Npn \__hit_cover_gauge:n #1
  {
    \vbox_set_top:Nn \l__hit_cover_box
      {
        \dim_set_eq:NN \hsize \l__hit_cover_width_dim
        \dim_set_eq:NN \linewidth \l__hit_cover_width_dim
        #1
      }
  }
\cs_new_protected:Npn \__hit_cover_gauge:Nn #1#2
  {
    \__hit_cover_gauge:n {#2}
    \dim_set:Nn #1 { \box_ht_plus_dp:N \l__hit_cover_box }
  }
\cs_new_protected:Npn \__hit_cover_gauge:cn #1#2
  { \exp_args:Nc \__hit_cover_gauge:Nn {#1} {#2} }
%    \end{macrocode}
%
% 各块的高只量一次，之后全是算术。相邻块的间距在这些切点上都是零，所以高度可加。
%    \begin{macrocode}
\cs_new_protected:Npn \__hit_cover_measure:n #1
  {
    \dim_set:Nn \l__hit_cover_top_dim
      { 1in + \voffset + \topmargin + \headheight + \headsep }
    \dim_set_eq:NN \l__hit_cover_width_dim \textwidth
    \dim_set_eq:NN \l__hit_cover_height_dim \textheight
    \__hit_cover_gauge:Nn \l__hit_cover_bc_dim
      { \hit@cover@blank@zh:n { 1 } }
    \__hit_cover_gauge:Nn \l__hit_cover_be_dim
      { \hit@cover@blank@en:n { 1 } }
    \int_step_inline:nnn { 0 } { 3 }
      {
        \__hit_cover_gauge:cn { l__hit_cover_ba_##1_dim }
          { \hit@cover@fixed@author:n {##1} }
        \__hit_cover_gauge:cn { l__hit_cover_bi_##1_dim }
          { \hit@cover@fixed@inst:n {##1} }
      }
    \__hit_cover_gauge:n { \hit@cover@author@line: }
    \dim_set:Nn \l__hit_cover_ha_dim { \box_ht_plus_dp:N \l__hit_cover_box }
    \__hit_cover_gauge:n { \hit@cover@inst@lines: }
    \dim_set:Nn \l__hit_cover_iasc_dim { \box_ht:N \l__hit_cover_box }
    \dim_set:Nn \l__hit_cover_hi_dim { \box_ht_plus_dp:N \l__hit_cover_box }
    \int_step_inline:nnn { 0 } { 2 }
      {
        \__hit_cover_gauge:cn { l__hit_cover_tu_##1_dim }
          { \hit@cover@upto@en:nn {#1} {##1} }
      }
  }
%    \end{macrocode}
%
% 一组取值下的几个关键量。都以磅计，存进浮点。
%
% 姓名到年月之间的跨度可能是零：最后一档让步把姓名与年月的固定空行都去掉，
% 这一组又取零个空行，而姓名那一块本身量出来是零高（v3.1e 的英文封面不排
% 中文姓名）。这时比例无从谈起，记零，得分只剩校名那一项。
%    \begin{macrocode}
\fp_new:N \l__hit_cover_inst_fp
\fp_new:N \l__hit_cover_span_fp
\fp_new:N \l__hit_cover_ratio_fp
\fp_new:N \l__hit_cover_total_fp
\fp_new:N \l__hit_cover_ref_fp
\fp_new:N \l__hit_cover_score_fp
\fp_new:N \l__hit_cover_best_fp
\int_new:N \l__hit_cover_near_int
\int_new:N \l__hit_cover_best_near_int
\cs_new_protected:Npn \__hit_cover_solve:nnnnn #1#2#3#4#5
  {
    \fp_set:Nn \l__hit_cover_span_fp
      {
        \dim_to_fp:n { \l__hit_cover_top_dim }
        + \dim_to_fp:n { \use:c { l__hit_cover_tu_#1_dim } }
      }
    \fp_set:Nn \l__hit_cover_ratio_fp
      {
        \l__hit_cover_span_fp
        + #2 * \dim_to_fp:n { \l__hit_cover_bc_dim }
        + \dim_to_fp:n { \use:c { l__hit_cover_ba_#5_dim } }
      }
    \fp_set:Nn \l__hit_cover_inst_fp
      {
        \l__hit_cover_ratio_fp
        + \dim_to_fp:n { \l__hit_cover_ha_dim }
        + #3 * \dim_to_fp:n { \l__hit_cover_bc_dim }
        + \dim_to_fp:n { \use:c { l__hit_cover_bi_#5_dim } }
        + #4 * \dim_to_fp:n { \l__hit_cover_be_dim }
        + \dim_to_fp:n { \l__hit_cover_iasc_dim }
      }
    \fp_set:Nn \l__hit_cover_total_fp
      {
        \l__hit_cover_inst_fp - \dim_to_fp:n { \l__hit_cover_iasc_dim }
        - \dim_to_fp:n { \l__hit_cover_top_dim }
        + \dim_to_fp:n { \l__hit_cover_hi_dim }
      }
    \fp_set:Nn \l__hit_cover_span_fp
      {
        \l__hit_cover_inst_fp - \dim_to_fp:n { \l__hit_cover_iasc_dim }
        - \l__hit_cover_span_fp
      }
    \fp_compare:nNnTF { \l__hit_cover_span_fp } = { 0 }
      { \fp_zero:N \l__hit_cover_ratio_fp }
      {
        \fp_set:Nn \l__hit_cover_ratio_fp
          {
            ( \l__hit_cover_ratio_fp
              + 0.5 * \dim_to_fp:n { \l__hit_cover_ha_dim }
              - \dim_to_fp:n { \l__hit_cover_top_dim }
              - \dim_to_fp:n { \use:c { l__hit_cover_tu_#1_dim } } )
            / \l__hit_cover_span_fp
          }
      }
  }
%    \end{macrocode}
%
% 姓名该落在什么比例上——不拍数，从范例自己算。Word 范例那一页用的是 2 / 2 / 6
% 这组空行数，把它代进去算出来的比例就是范例的比例。题目变长变短时按这个比例
% 去凑，姓名的视觉位置才不变。
%    \begin{macrocode}
\dim_new:N \l__hit_cover_keep_dim
\cs_new_protected:Npn \__hit_cover_ref_ratio:
  {
    \dim_set_eq:Nc \l__hit_cover_keep_dim { l__hit_cover_tu_0_dim }
    \dim_zero:c { l__hit_cover_tu_0_dim }
    \__hit_cover_solve:nnnnn { 0 } { 2 } { 2 } { 6 } { 0 }
    \fp_set_eq:NN \l__hit_cover_ref_fp \l__hit_cover_ratio_fp
    \dim_set_eq:cN { l__hit_cover_tu_0_dim } \l__hit_cover_keep_dim
  }
%    \end{macrocode}
%
% 挑一组空行数。目标按优先级：一、校名与年月的位置尽可能不变；二、姓名在视觉上
% 落在范例里那个比例上。两项都是长度，加权相加：
% \[ 得分 = 3 \times |校名偏差| + |姓名偏离该在的位置| \]
% 校名计三倍权，也就是宁可姓名偏三磅，也不让校名偏一磅。不用字典序：校名的可达
% 落点很密，字典序下第一项几乎总能压到零点几磅，第二项就成了摆设。平局看谁离
% Word 自己那组取值（2 / 2 / 6）近。
%    \begin{macrocode}
\int_new:N \l__hit_cover_na_int
\int_new:N \l__hit_cover_nb_int
\int_new:N \l__hit_cover_nc_int
\bool_new:N \l__hit_cover_found_bool
\int_new:N \l__hit_cover_ia_int
\int_new:N \l__hit_cover_ib_int
\int_new:N \l__hit_cover_ic_int
\cs_new_protected:Npn \__hit_cover_best:nn #1#2
  {
    \bool_set_false:N \l__hit_cover_found_bool
    \int_step_inline:nnn { 0 } { 272 }
      {
        \int_set:Nn \l__hit_cover_ia_int { \int_div_truncate:nn {##1} { 39 } }
        \int_set:Nn \l__hit_cover_ib_int
          { \int_div_truncate:nn { \int_mod:nn {##1} { 39 } } { 13 } }
        \int_set:Nn \l__hit_cover_ic_int { \int_mod:nn {##1} { 13 } }
        \__hit_cover_try:nn {#1} {#2}
      }
  }
\cs_new_protected:Npn \__hit_cover_try:nn #1#2
  {
    \__hit_cover_solve:nnnnn {#1}
      { \l__hit_cover_ia_int } { \l__hit_cover_ib_int }
      { \l__hit_cover_ic_int } {#2}
    \fp_compare:nNnT { \l__hit_cover_total_fp } <
      { \dim_to_fp:n { \l__hit_cover_height_dim } }
      {
        \fp_set:Nn \l__hit_cover_score_fp
          {
            3 * abs ( \l__hit_cover_inst_fp
              - \dim_to_fp:n { \l__hit_cover_inst_target_dim } )
            + abs ( \l__hit_cover_ratio_fp - \l__hit_cover_ref_fp )
              * \l__hit_cover_span_fp
          }
        \int_set:Nn \l__hit_cover_near_int
          {
            \int_abs:n { \l__hit_cover_ia_int - 2 }
            + \int_abs:n { \l__hit_cover_ib_int - 2 }
            + \int_abs:n { \l__hit_cover_ic_int - 6 }
          }
        \bool_if:nT
          {
            ! \l__hit_cover_found_bool
            || \fp_compare_p:nNn { \l__hit_cover_score_fp } <
                 { \l__hit_cover_best_fp }
            || (
                 \fp_compare_p:nNn { \l__hit_cover_score_fp } =
                   { \l__hit_cover_best_fp }
                 && \int_compare_p:nNn { \l__hit_cover_near_int } <
                      { \l__hit_cover_best_near_int }
               )
          }
          {
            \bool_set_true:N \l__hit_cover_found_bool
            \fp_set_eq:NN \l__hit_cover_best_fp \l__hit_cover_score_fp
            \int_set_eq:NN \l__hit_cover_best_near_int \l__hit_cover_near_int
            \int_set_eq:NN \l__hit_cover_na_int \l__hit_cover_ia_int
            \int_set_eq:NN \l__hit_cover_nb_int \l__hit_cover_ib_int
            \int_set_eq:NN \l__hit_cover_nc_int \l__hit_cover_ic_int
          }
      }
  }
%    \end{macrocode}
%
% 结构性让步的六档。挤占中英题目之间的空行、把那两段 1.25 倍改成单倍、再到把
% 它们删掉——每一档都在动 Word 的段落结构，代价依次变大。第一项是挤占数，
% 第二项是让步级。
%    \begin{macrocode}
\dim_const:Nn \c__hit_cover_threshold_dim { 10bp }
%    \end{macrocode}
%
% \begin{macro}{\hit@cover@not@postdoc}
% 本硕博共用的封面一。先换上封面这一节自己的版面，再量、挑一组空行数、
% 照挑出来的那组排，排完把版面还回去。
%
% Word 的分节符让封面那一节有自己的页面设置。与正文只差一处：\texttt{docGrid}
% 的 \texttt{linePitch} 是 395 缇（19.75\,bp），正文是 383（19.15\,bp）。封面上
% 只有中文题目那一段贴网格，所以这一处只落在它身上：自然行高 1.296875×22
% = 28.53 装不下一格，占两格居中，格子大 0.6 则行上余量也大 0.6，基线跟着低 0.6。
%
% 字号按塞不塞得下定，不按排得好不好定。规范给的是二号，「题目太长时可用
% 小二号」——降号是规范明写允许的一档，但也只在二号排不下时才用，不为了排得
% 好看去缩号。所以字号是让步阶梯的最外层。旧版按宽度判（超过 250mm 就降小二），
% 那等于英文题目一过两行就降号，可封面上排三四行英文完全放得下。
%
% 结构档只在排得不够好时才进来。不动结构就能排到阈值以内就不动——结构让步的
% 本意是不为小便宜动 Word 的段落结构。
%    \begin{macrocode}
\fp_new:N \l__hit_cover_gbest_fp
\int_new:N \l__hit_cover_gna_int
\int_new:N \l__hit_cover_gnb_int
\int_new:N \l__hit_cover_gnc_int
\int_new:N \l__hit_cover_gsq_int
\int_new:N \l__hit_cover_grelax_int
\bool_new:N \l__hit_cover_gfound_bool
\tl_new:N \l__hit_cover_size_tl
\cs_new_protected:Npn \__hit_cover_adopt:nn #1#2
  {
    \bool_set_true:N \l__hit_cover_gfound_bool
    \fp_set_eq:NN \l__hit_cover_gbest_fp \l__hit_cover_best_fp
    \int_set_eq:NN \l__hit_cover_gna_int \l__hit_cover_na_int
    \int_set_eq:NN \l__hit_cover_gnb_int \l__hit_cover_nb_int
    \int_set_eq:NN \l__hit_cover_gnc_int \l__hit_cover_nc_int
    \int_set:Nn \l__hit_cover_gsq_int {#1}
    \int_set:Nn \l__hit_cover_grelax_int {#2}
  }
\cs_new_protected:Npn \__hit_cover_pick:
  {
    \bool_set_false:N \l__hit_cover_gfound_bool
    \str_if_eq:VnTF \g_hit_cover_title_en_xiaoer_tl { true }
      { \tl_set:Nn \l__hit_cover_size_tl { xiaoer } }
      { \tl_set:Nn \l__hit_cover_size_tl { erhao } }
    \exp_args:NV \__hit_cover_measure:n \l__hit_cover_size_tl
    \__hit_cover_ref_ratio:
    \__hit_cover_best:nn { 0 } { 0 }
    \bool_lazy_and:nnT
      { ! \bool_if_p:N \l__hit_cover_found_bool }
      { \str_if_eq_p:Vn \g_hit_cover_title_en_xiaoer_tl { auto } }
      {
        \tl_set:Nn \l__hit_cover_size_tl { xiaoer }
        \__hit_cover_measure:n { xiaoer }
        \__hit_cover_best:nn { 0 } { 0 }
      }
    \bool_if:nTF
      {
        \l__hit_cover_found_bool
        && \fp_compare_p:nNn { \l__hit_cover_best_fp } <
             { \dim_to_fp:n { \c__hit_cover_threshold_dim } }
      }
      { \__hit_cover_adopt:nn { 0 } { 0 } }
      {
        \__hit_cover_rung:nn { 0 } { 0 }
        \__hit_cover_rung:nn { 1 } { 0 }
        \__hit_cover_rung:nn { 2 } { 0 }
        \__hit_cover_rung:nn { 2 } { 1 }
        \__hit_cover_rung:nn { 2 } { 2 }
        \__hit_cover_rung:nn { 2 } { 3 }
      }
  }
\cs_new_protected:Npn \__hit_cover_rung:nn #1#2
  {
    \__hit_cover_best:nn {#1} {#2}
    \bool_if:nT
      {
        \l__hit_cover_found_bool
        && (
             ! \l__hit_cover_gfound_bool
             || \fp_compare_p:nNn { \l__hit_cover_best_fp } <
                  { \l__hit_cover_gbest_fp }
           )
      }
      { \__hit_cover_adopt:nn {#1} {#2} }
  }
%    \end{macrocode}
%
% 让步全用尽还是装不下时不拦编译，按压得最紧的一组照排，另发一条警告。为了一页
% 封面卡住整篇文档不合适。
%    \begin{macrocode}
\cs_new_protected:Npn \hit@cover@not@postdoc
  {
    \hit@layout@push@part:n { cover }
    \group_begin:
      \dim_zero:N \parindent
      \dim_zero:N \parskip
      \bool_if:NTF \g__hit_bachelor_bool
        { \dim_set:Nn \l__hit_cover_inst_target_dim { 634.46bp } }
        { \dim_set:Nn \l__hit_cover_inst_target_dim { 689.36bp } }
      \__hit_cover_pick:
      \bool_if:NF \l__hit_cover_gfound_bool
        {
          \msg_warning:nn { hithesis } { cover-overflow }
          \tl_set:Nn \l__hit_cover_size_tl { xiaoer }
          \int_set:Nn \l__hit_cover_gsq_int { 2 }
          \int_zero:N \l__hit_cover_gna_int
          \int_zero:N \l__hit_cover_gnb_int
          \int_zero:N \l__hit_cover_gnc_int
          \int_set:Nn \l__hit_cover_grelax_int { 3 }
        }
      \exp_args:NV \hit@cover@upto@en:nn \l__hit_cover_size_tl
        { \l__hit_cover_gsq_int }
      \hit@cover@blank@zh:n { \l__hit_cover_gna_int }
      \hit@cover@fixed@author:n { \l__hit_cover_grelax_int }
      \hit@cover@author@line:
      \hit@cover@blank@zh:n { \l__hit_cover_gnb_int }
      \hit@cover@fixed@inst:n { \l__hit_cover_grelax_int }
      \hit@cover@blank@en:n { \l__hit_cover_gnc_int }
      \hit@cover@inst@lines:
    \group_end:
    \hit@layout@pop@part:
  }
%    \end{macrocode}
%    \end{macro}

%    \begin{macro}{\makecover}
%    \begin{macrocode}

%    \end{macrocode}
%
% \changes{v2.0.10}{2019/06/25}{此处添加提交图书馆电子版的逻辑}
% \changes{v3.0.0}{2020/05/21}{博后没有英文封皮}
% 封面分派。前两张封面人人都排；博后与本科没有英文封皮。翻页动作取决于
% \option{style/library}，交图书馆电子版时用 \cs{clearpage}，否则 \cs{cleardoublepage}。
%    \begin{macrocode}
\cs_new_protected:Npn \hit@clear@page
  { \exp_args:NV \hit@openright:n \g__hit_matter_now_tl }
\cs_new_protected:Npn \__hit_cover_break:
  { \hit@openright:n { cover } }
%    \end{macrocode}
%
% \emph{排过就记一笔。} \file{hit-structure.dtx} 里有个兜底：文末发现没调过
% \cs{hitfrontmatter} 就补调一次，防的是「用户忘了写骨架、整份文档只有正文」。
% 可 v3.1e 那种老写法用的是 \cs{frontmatter} 加 \cs{makecover}，不走骨架命令，
% 标志位没置上，兜底就在文末把封面扉页摘要又排了一遍——旧版面兼容回归里实测
% 多出 8 页，正是这一份。\cs{makecover} 排的就是前置，排完置位，兜底自然不发。
%    \begin{macrocode}
\cs_set_protected:Npn \makecover
  {
    \bool_gset_true:N \g__hit_matter_front_bool
    \hit@write@pdf@metadata
    \phantomsection
    \pdfbookmark [0] { \hit@info@title@cover@zh } { ctitle }  % ^^A PDF 目标名，保持稳定
    \hit@fontsize@x{xiaosi}{1}
    \begin { titlepage }
%    \end{macrocode}
%
% 封面上的宽度大多写成 \cs{ccwd} 的倍数（\cs{hspace}\marg{7\cs{ccwd}}、
% \cs{makebox}\oarg{6\cs{ccwd}}），而 \cs{frontmatter} 里的 \cs{ziju} 让 \cs{ccwd}
% 比一个字宽大 5.6\%，凑出来的表格因此会超出版心几个点。超出去的是空列的留白，
% 版面上什么都没被裁，但每份博后文档要为此报一条 Overfull。这里把容差放宽到
% 一个字宽，纯粹是不报，排版一点不动。
%    \begin{macrocode}
      \dim_set:Nn \hfuzz { 1em }
      \hit@cover
      \__hit_cover_break:
      \hit@titlepage
      \__hit_cover_break:
      \bool_lazy_or:nnF
        { \bool_if_p:N \g__hit_postdoc_bool }
        { \bool_if_p:N \g__hit_bachelor_bool }
        {
          \phantomsection
          \pdfbookmark [0] { \hit@info@title@cover@en } { etitle }  % ^^A 同上
          \hit@titlepage@en
          \__hit_cover_break:
        }
    \end { titlepage }
    \normalsize
    \hit@make@abstract
  }
%    \end{macrocode}
%    \end{macro}
%
%    \begin{macrocode}
  }
%</final-cover>
%    \end{macrocode}
%
% \Finale
